\chapter{Playwright}

\section{Ce este Playwright?}

Playwright este un framework open-source de automatizare a browserelor, dezvoltat de Microsoft
și lansat în 2020. Este conceput pentru a testa aplicații web moderne, oferind suport nativ
pentru trei motoare de browser: \textbf{Chromium}, \textbf{Firefox} și \textbf{WebKit} (Safari).
Spre deosebire de predecesori precum Selenium, Playwright controlează browserul printr-un
protocol binar propriu (DevTools Protocol), ceea ce îi conferă viteză și fiabilitate superioare.

Playwright pune la dispoziție API-uri pentru Python, JavaScript/TypeScript, Java și C\#.
În cadrul acestui proiect a fost utilizat \textbf{API-ul sincron Python}, integrat cu
runner-ul de teste \texttt{pytest}.

\section{Facilități principale}

\subsection{Suport multi-browser}

Playwright rulează aceleași teste pe Chromium, Firefox și WebKit fără modificări de cod.
Fiecare motor este distribuit împreună cu biblioteca, eliminând dependența de versiunile
instalate local pe mașina de test.

\subsection{Auto-wait}

Playwright implementează un mecanism de așteptare automată (\textit{auto-wait}): înainte de
a interacționa cu un element, framework-ul verifică automat că acesta este vizibil, activat
și stabil în DOM. Acest comportament elimină necesitatea apelurilor explicite de tip
\texttt{sleep()} sau \texttt{wait\_for\_element()} în marea majoritate a cazurilor.

\subsection{Generare automată de teste --- Codegen}

Playwright include un instrument de înregistrare a interacțiunilor numit \texttt{codegen}.
La rularea comenzii, se deschide un browser în care utilizatorul interacționează normal
cu aplicația, iar Playwright generează automat codul de test corespunzător în timp real.

\begin{lstlisting}[language=bash]
playwright codegen http://localhost:8080
\end{lstlisting}

Codul generat poate fi copiat direct în fișierul de test și rafinat ulterior. Aceasta
accelerează semnificativ scrierea testelor pentru fluxuri noi, eliminând nevoia de a
identifica manual selectorii elementelor.

\subsection{Interceptarea rețelei}

API-ul \texttt{page.expect\_response()} permite atașarea unui context de așteptare la un
apel de rețea specific, identificat prin URL sau expresie regulată. Testul este blocat până
când răspunsul este recepționat, moment în care se poate inspecta statusul sau corpul
răspunsului. Aceasta este o alternativă robustă la așteptările bazate pe timp fix.

\subsection{Locatori rezilienți}

Playwright introduce conceptul de \textit{Locator}, o referință lazy către un element din DOM.
Spre deosebire de un \texttt{ElementHandle} clasic (capturat la un moment de timp), un Locator
re-caută elementul la fiecare interacțiune, adaptându-se la re-randările dinamice ale paginii.

\subsection{Izolarea contextelor de browser}

Fiecare test poate rula într-un context de browser separat (\texttt{BrowserContext}), echivalent
cu un profil de utilizator distinct. Contextele sunt izolate: cookie-uri, localStorage și cache-ul
nu se propagă între ele. Plugin-ul \texttt{pytest-playwright} creează automat câte un context și
o pagină pentru fiecare funcție de test.

\subsection{Mod Trace și raportare}

Playwright include un instrument de înregistrare a execuției (\textit{Trace Viewer}) care
captează screenshot-uri, snapshot-uri DOM, apeluri de rețea și timeline-ul acțiunilor.
Fișierul de trace poate fi inspectat ulterior în browser, fără a rula din nou testul.

\section{Configurarea proiectului}

\subsection{Dependențe}

Dependențele proiectului sunt definite în \texttt{requirements.txt}:

\begin{lstlisting}[language=Python, caption={requirements.txt}]
playwright==1.58.0
pytest==9.0.3
pytest-playwright==0.7.2
pytest-base-url==2.1.0
playwright-stealth==2.0.3
\end{lstlisting}

Instalarea se face cu:

\begin{lstlisting}[language=bash]
pip install -r requirements.txt
playwright install
\end{lstlisting}

Comanda \texttt{playwright install} descarcă binarele pentru Chromium, Firefox și WebKit.

\subsection{Bypass Cloudflare cu playwright-stealth}

Instanța publică \texttt{demo.nopcommerce.com} este protejată de Cloudflare, care detectează
și blochează sesiunile de browser automatizate prin analiza fingerprint-ului JavaScript
(prezența variabilei \texttt{navigator.webdriver}, absența plugin-urilor, timingul evenimentelor).

Pentru a ocoli această protecție, proiectul utilizează biblioteca \texttt{playwright-stealth},
care modifică proprietățile expuse de browser astfel încât sesiunea automată să fie
indistinctibilă de una umană. Integrarea se face printr-un fixture \texttt{pytest} cu
\texttt{autouse=True}, aplicat global tuturor testelor:

\begin{lstlisting}[language=Python, caption={Fixture global pentru stealth mode}]
@pytest.fixture(autouse=True)
def apply_stealth(page: Page):
    Stealth().apply_stealth_sync(page)
\end{lstlisting}

Prin declararea \texttt{autouse=True}, fixture-ul este injectat automat în fiecare test
fără a fi necesar să fie menționat explicit în semnătura funcției.

\section{Implementarea testelor}

\subsection{Structura unui test}

Fiecare test este o funcție Python prefixată cu \texttt{test\_}, care primește ca argument
obiectul \texttt{page} furnizat de plugin-ul \texttt{pytest-playwright}. Acesta reprezintă
o instanță de pagină browser gata de utilizare, izolată față de celelalte teste.

\begin{lstlisting}[language=Python, caption={Structura de bază a unui test Playwright}]
def test_guest_shopping(page: Page):
    page.goto(BASE_URL)
    page.locator("a.menu__link[href='/computers']").hover()
    page.locator("a[href='/notebooks']").click()
    # ...
    expect(page.locator(".order-completed .title")).to_have_text(
        "Your order has been successfully processed!"
    )
\end{lstlisting}

\subsection{Interceptarea cererilor AJAX}

Un pattern central în suita de teste este așteptarea răspunsurilor AJAX înainte de a
continua execuția. Fără această sincronizare, testul ar interacționa cu elemente care
nu reflectă încă rezultatul operației asincrone.

Exemplul următor verifică că filtrarea produselor după RAM s-a încheiat cu succes înainte
de a continua:

\begin{lstlisting}[language=Python, caption={Interceptare răspuns AJAX la aplicarea filtrului}]
with page.expect_response(
    lambda response: "demo.nopcommerce.com" in response.url
    and response.request.resource_type in ["xhr", "fetch"]
) as response_info:
    page.locator("input[id='attribute-option-9']").check()

assert response_info.value.status == 200, \
    "Filtrarea asincronă (AJAX) nu a returnat status 200"
\end{lstlisting}

Blocul \texttt{with page.expect\_response(...)} funcționează ca un context manager: acțiunea
din interior declanșează cererea, iar Playwright așteaptă răspunsul înainte de a ieși din bloc.

\subsection{Verificarea prețurilor dinamice}

Workflow-ul 2 (Configuratorul PC) implică selectarea unor componente care modifică prețul
afișat prin AJAX. Testul extrage prețul înainte și după selecție și verifică matematic că
actualizarea a avut loc:

\begin{lstlisting}[language=Python, caption={Verificare calcul preț dinamic}]
base_price_text = page.locator("div.product-price span").inner_text()
base_price = float(re.sub(r'[^\d.]', '', base_price_text))

with page.expect_response(re.compile(r".*productdetails_attributechange.*",
                                     re.IGNORECASE)):
    page.locator("input[name='product_attribute_3'][value='7']").check()

new_price_text = page.locator("div.product-price span").inner_text()
new_price = float(re.sub(r'[^\d.]', '', new_price_text))

assert new_price > base_price, \
    f"Pretul nu s-a actualizat corect! {new_price} vs {base_price}"
\end{lstlisting}

\subsection{Soft assertions}

Un \textit{soft assert} este o verificare al cărei eșec nu oprește execuția testului.
Playwright nu oferă un mecanism nativ pentru soft assertions, însă acestea se implementează
simplu prin capturarea excepției:

\begin{lstlisting}[language=Python, caption={Implementare soft assert pentru cuponul invalid}]
try:
    page.wait_for_selector("div.message-failure", timeout=5000)
    expect(page.locator("div.message-failure").first).to_have_text(
        re.compile(r"coupon", re.IGNORECASE)
    )
except (AssertionError, Exception) as e:
    print(f"\nMesajul de eroare pentru cupon nu a aparut. "
          f"Soft-assert ignorat. Detalii: {e}")
    pass
\end{lstlisting}

Dacă mesajul de eroare nu apare (comportament întâlnit uneori pe instanța demo),
testul înregistrează situația în consolă și continuă spre pașii următori.

\subsection{Gestionarea sesiunilor persistente}

Workflow-ul 6 testează că datele (coș, adrese) supraviețuiesc unui ciclu de logout/login.
Pentru aceasta, testul utilizează un cont cu credențiale fixe, creat o singură dată:

\begin{lstlisting}[language=Python, caption={Funcție helper pentru autentificare}]
PERSISTENT_EMAIL = "persistent.tester2026@example.com"
PERSISTENT_PASSWORD = "PersistTest123!"

def _login(page: Page):
    page.goto(f"{BASE_URL}/login")
    page.locator("#Email").fill(PERSISTENT_EMAIL)
    page.locator("#Password").fill(PERSISTENT_PASSWORD)
    page.locator("button.login-button").click()
    page.wait_for_load_state("networkidle")
\end{lstlisting}

Funcția \texttt{\_ensure\_account\_exists()} încearcă înregistrarea contului; dacă
primește o eroare (contul există deja), continuă direct cu autentificarea ---
eliminând necesitatea unui setup manual înainte de prima rulare.

\subsection{Adresă de email unică per rulare}

Workflow-ul 5 (înregistrare cont) necesită un email unic la fiecare rulare pentru a evita
eroarea \textit{``account already exists''}. Soluția este generarea unui timestamp:

\begin{lstlisting}[language=Python, caption={Generare email unic cu timestamp}]
unique_email = f"testuser_{int(time.time())}@example.com"
\end{lstlisting}

\section{Avantaje}

\begin{itemize}
    \item \textbf{Suport multi-browser nativ}: Chromium, Firefox și WebKit testate cu același cod,
          fără configurare suplimentară.
    \item \textbf{Interceptare rețea puternică}: \texttt{expect\_response()} și
          \texttt{route()} permit sincronizare și modificare a cererilor HTTP direct din test.
    \item \textbf{Auto-wait robust}: elimină majoritatea așteptărilor explicite, reducând
          fragilitatea testelor la variații de latență.
    \item \textbf{Izolare contexte}: fiecare test pornește cu un browser curat, fără interferențe
          între sesiuni.
    \item \textbf{Ecosistem Python}: integrare naturală cu \texttt{pytest}, fixtures,
          parametrizare, și librăriile de analiză/date din ecosistemul Python.
    \item \textbf{Bypass Cloudflare}: prin \texttt{playwright-stealth}, testele pot rula
          pe instanțe protejate fără a necesita un mediu local.
    \item \textbf{Trace Viewer}: debugging post-mortem fără re-rularea testelor.
\end{itemize}

\section{Dezavantaje}

\begin{itemize}
    \item \textbf{Configurare inițială mai complexă}: necesită instalarea separată a binarelor
          de browser și configurarea mediului Python, față de un simplu \texttt{npm install}.
    \item \textbf{Verbozitate}: codul Python tinde să fie mai explicit față de echivalentul
          Cypress, mai ales în cazul selectoarelor și așteptărilor condiționale.
    \item \textbf{Fără Time Travel nativ}: Playwright nu oferă un echivalent al funcției
          de \textit{time travel debugging} din Cypress App --- reluarea pas cu pas a
          execuției necesită activarea Trace Viewer înainte de rulare.
    \item \textbf{Soft assertions manuale}: absența unui mecanism nativ pentru soft assertions
          obligă la construcții \texttt{try/except} repetitive.
    \item \textbf{Dependență de stealth pentru site-uri protejate}: pe instanțe publice
          cu protecție anti-bot, este necesară o bibliotecă externă care poate deveni
          incompatibilă cu versiuni noi de Playwright.
\end{itemize}
